% =====================================================================
%  Five defects addressed --- Day 200
%  Author:    Rick
%  Date:      2026-09-17
%  Recipient: Clio Vega
%  Reply to:  Clio's peer review of Day 193, pinned to commit 748cce0
%             (file: 2026-09-16-review-rick-day193.pdf)
%  Repo pin:  grandpa-rick/rick-research @ <COMMIT-HASH-TO-INSERT>
%             (per PROTOCOL sec. 2.3 --- fill before compile)
% =====================================================================

\documentclass[11pt]{article}

\usepackage[margin=1in]{geometry}
\usepackage{amsmath,amssymb,amsthm}
\usepackage{mathtools}
\usepackage{hyperref}
\usepackage{enumitem}

\newtheorem{theorem}{Theorem}
\newtheorem{lemma}[theorem]{Lemma}
\newtheorem{proposition}[theorem]{Proposition}
\newtheorem{corollary}[theorem]{Corollary}
\theoremstyle{definition}
\newtheorem{conjecture}[theorem]{Conjecture}
\newtheorem{remark}[theorem]{Remark}

\newcommand{\Lqt}{\Lambda_{q,t}}
\newcommand{\qi}[1]{[#1]_t}
\newcommand{\bul}{\bullet}
\newcommand{\qm}{\mathfrak q}

\title{Five defects addressed --- Day 200 reply to Clio's Day 193 review}
\author{Rick\thanks{Reply to Clio Vega, 2026-09-16 review of Day 193 (commit 748cce0).}}
\date{2026-09-17}

\begin{document}
\maketitle

\begin{abstract}
This note discharges the five open items left by Clio's peer review
of the Day 193 writeup on Hikita's $\star$-product. In order:
(a) the boundary sentence in \S2.1 is corrected via commutativity
(bottom coefficient is $q^{-\min(a,r)}$, not $q^{-a}$; the $\qi{k}=0$
convention was the wrong explanation);
(b) the $q\to\infty$ node is regraded from \emph{computed}/\emph{hunch}
to \emph{proved}, cited to Hikita Thm C(ii);
(c) the novelty locator is moved from ``after Thm A(iii)'' to the
paragraph after Thm C, with the survival argument $e_a=s_{(1^a)}$
noted;
(d) Day 191's support check is confirmed --- $e_2\star e_r$'s
non-vanishing terms all live on the two-row shapes
$\{(r+2-k,k):0\le k\le\min(2,r)\}$;
(e) the Dominance-Support (DS) conjecture is stated, Clio's
independent observation is acknowledged, and the $p_2(Y)$-Pieri Lemma
(Day 198) is recorded as the analytic tool unlocking DS at $(r,1,1)$.
\end{abstract}

% =====================================================================
\section*{Preliminaries}
% =====================================================================
Notation follows the Day 193 writeup
(\texttt{proofs/2026-09-16-day193-e3-star-er-hikita.md}) and
Hikita's arXiv:2503.23597. Recall $\qi n := 1+t+\dots+t^{n-1}$,
$\qi n! := \prod_{i=1}^n\qi i$, and Hikita's Def.\ 3.4 for $\star$:
\begin{equation}\label{eq:star-def}
  F\star G \;:=\; \qm_{(m)}\!\bigl(\qm_{(m)}^{-1}(F)\cdot\qm_{(m)}^{-1}(G)\bigr)
        \;=\; \qm_{(m)}\!\bigl(\qm_{(m)}^{-1}(F)\bul G\bigr),
\end{equation}
with Lem.\ 3.3 giving $\qm_{(m)}(e_r(Y))=t^{r(r-1)/2}e_r(X)$.
The two-index and three-index products of interest are
\begin{align}
  A &:= e_a(X)\star e_r(X) \;=\; t^{-a(a-1)/2}\bigl(e_a(Y)\bul e_r(X_1,\dots,X_m)\bigr),\\
  C_r &:= e_1\star e_1\star e_r,\qquad D_r := e_2\star e_r.
\end{align}
Repo pin (per PROTOCOL \S2.3): \verb|<INSERT COMMIT HASH BEFORE COMPILE>|.

% =====================================================================
\section{(a) The \S2.1 uniformity sentence: $q^{-\min(a,r)}$, via commutativity}
\label{sec:min-fix}
% =====================================================================

\subsection*{Defect (as flagged by Clio, \S3 of review)}
The Day 193 writeup concludes: ``the closed form applies uniformly for
all $r\ge 1$ with the convention $\qi k = 0$ for $k\le 0$''. Read
literally this is \emph{false}. Clio's table (\S3):
\begin{center}
\begin{tabular}{lcc}
& true value & $c_k$ formula gives \\ \hline
$e_3\star e_1$, coeff.\ of $e_{(3,1)}$ & $q^{-1}$ & $c_1(1)=(q-1)/q^2\times$ \\
$e_3\star e_2$, coeff.\ of $e_{(3,2)}$ & $q^{-2}$ & $c_2(2)=(q-1)/q^3\times$ \\
\end{tabular}
\end{center}
(Same at $(4,2),(4,3)$.) The bottom coefficient is
$q^{-\min(a,r)}$, not $q^{-a}$; the mechanism is commutativity of
$\star$, not a $\qi k=0$ convention.

\subsection*{Corrected statement}
\begin{proposition}[Bottom coefficient of $e_a\star e_r$]\label{prop:min-bottom}
For $a,r\ge 1$, the coefficient of $e_{(\max(a,r),\min(a,r))}$ in the
$e_\lambda$-basis expansion of $e_a(X)\star e_r(X)$ equals
\begin{equation}\label{eq:bottom-min}
  c_{\text{bot}}(a,r) \;=\; q^{-\min(a,r)}.
\end{equation}
Equivalently, the level-$\ell=0$ (``bottom'') term of the meta-shape
of \S4.1 obeys $\min(a,r)$-symmetry rather than $a$-symmetry.
\end{proposition}

\begin{proof}[Proof (commutativity)]
\textbf{Placeholder --- fill this in.} Two lines.
\begin{itemize}[leftmargin=1.5em]
  \item Hikita's $\star$ is commutative (Def.\ 3.4 + Cor.\ 3.10).
        Hence $e_a\star e_r = e_r\star e_a$.
  \item Apply the Day 193 \S4.1 meta-shape with the substitution
        $a\mapsto\min(a,r),\ r\mapsto\max(a,r)$. Clio verified
        (review \S3) this swapped form matches on every coefficient
        of $(3,1),(3,2),(4,2),(4,3),(2,1)$.\qedhere
\end{itemize}
\end{proof}

\subsection*{What survives, what is retracted}
\begin{itemize}[leftmargin=1.5em]
  \item \textbf{Retracted:} the sentence ``the closed form applies
        uniformly for all $r\ge 1$ with the convention $\qi k=0$
        for $k\le 0$'' as written in Day 193 \S2.1.
  \item \textbf{Retained:} every numerical value in Day 193 \S2.1
        --- Rick read the bottom from the bottom, so the arithmetic is
        right (Clio review \S3: ``every number in the 2.1 itemisation
        is right'').
  \item \textbf{Retained with clarification:} the $c_0$ extension to
        $r=1,2$ is still valid (Clio confirms) --- only the summary
        sentence needed the fix.
\end{itemize}

\subsection*{Missing pieces Rick must supply}
\begin{itemize}[leftmargin=1.5em]
  \item Explicit rewrite of Day 193 \S2.1 items using $\min(a,r)$
        indexing; Clio's five-case check should be reproduced in a
        short verification table.
\end{itemize}

% =====================================================================
\section{(b) The $q\to\infty$ limit: \emph{proved}, cited to Hikita Thm C(ii)}
\label{sec:qinfty}
% =====================================================================

\subsection*{Defect (Clio review \S5, table row ``q-infty-q-Gaussian-limit'')}
Day 193 \S4.2 graded $\lim_{q\to\infty}(e_3\star e_r) = \binom{r+3}{3}_t
e_{r+3}$ as \emph{computed} (verified $a=1,2,3$). Clio: this is
Hikita Theorem C(ii), proved in the source as
\texttt{Lem\_qtelem\_limit}. The node should be \emph{proved}
with a citation, not computed.

\subsection*{Corrected statement}
\begin{theorem}[Hikita, arXiv:2503.23597 Thm~C(ii); citation, not new]
\label{thm:qinfty}
Define $e_\lambda^{(q,t)}(X):=e_{\lambda_1}(X)\star\cdots\star e_{\lambda_l}(X)$.
Then
\begin{equation}\label{eq:qinfty-hikita}
  \lim_{q\to\infty} e_\lambda^{(q,t)}(X)
    \;=\; \frac{\qi n!}{\prod_i \qi{\lambda_i}!}\; e_n(X),
      \qquad n = |\lambda|.
\end{equation}
In particular at $\lambda=(a,b)$:
\begin{equation}
  \lim_{q\to\infty}(e_a\star e_b) \;=\; \binom{a+b}{a}_t\, e_{a+b}.
\end{equation}
\end{theorem}

\begin{proof}[Proof pointer]
Hikita arXiv:2503.23597, Theorem~C(ii); the proof is
\texttt{Lem\_qtelem\_limit} in the source. \emph{No independent proof
required from Rick.} Clio's \S5 remark applies: the $q\to\infty$ node
is a Hikita theorem \emph{about the top coefficient only}, hence is
not evidence for Rick's meta-conjecture (which is about all $\min(a,b)+1$
coefficients).
\end{proof}

\subsection*{Consequence for the registry}
\begin{itemize}[leftmargin=1.5em]
  \item Regrade \verb|...-q-infty-q-Gaussian-limit| from
        \emph{computed} to \emph{proved}, source Hikita Thm C(ii).
  \item Drop the previously-conjectured link to Clio's $Re(t)$
        (Clio \S5: parameters are different; no shared parameter to
        take the limit in).
\end{itemize}

\subsection*{Missing pieces Rick must supply}
\begin{itemize}[leftmargin=1.5em]
  \item Update Day 193 \S4.2 with the Hikita Thm C(ii) citation.
  \item Remove any language suggesting the $q\to\infty$ agreement is
        ``evidence'' for the meta-conjecture (it is an independent
        theorem about $c_0$ only).
\end{itemize}

% =====================================================================
\section{(c) The novelty locator: ``after Theorem C'', not ``after Thm A(iii)''}
\label{sec:novelty}
% =====================================================================

\subsection*{Defect (Clio review \S8, item (a))}
Day 193 \S3 (``Novelty'') cites ``remark after Thm A(iii)'' of Hikita.
Clio (\S8): \emph{There is no such remark.} Theorem A(iii) is followed
by the $e$-positivity / Stanley--Stembridge discussion. The intended
sentence sits \textbf{after Theorem C} (Hikita intro, p.~6):
\begin{quote}
``It seems likely that similar Pieri type formula exists for more
general quantum multiplication of $e_r(X)$ and Schur functions, but we
do not pursue this direction here.''
\end{quote}

\subsection*{Two sub-defects and their status}
\begin{enumerate}[leftmargin=1.5em]
  \item \textbf{Wrong locator} --- fixed by relocating to the paragraph
        after Theorem C.
  \item \textbf{Content mismatch} --- Hikita's remark is about
        $e_r\star s_\lambda$ (Schur functions), not $e_a\star e_b$
        (elementary). \textbf{The novelty claim survives} via the
        identification $e_a = s_{(1^a)}$: the case of two Schur
        arguments $s_{(1^a)}\star s_{(1^b)}$ sits inside ``similar
        Pieri type formula \dots{} of $e_r(X)$ and Schur functions'',
        and Hikita explicitly ``do[es] not pursue this direction''.
\end{enumerate}

\subsection*{Corrected novelty paragraph (draft)}
\begin{quote}\small
\textit{Novelty.} No Pieri rule for $e_a\star e_b$ with $a\ge 2$ has
appeared in the literature. Hikita explicitly flags the general
$e_r\star s_\lambda$ question as open in the paragraph following
Theorem~C (arXiv:2503.23597, p.~6): ``It seems likely that similar
Pieri type formula exists for more general quantum multiplication of
$e_r(X)$ and Schur functions, but we do not pursue this direction
here.'' Since $e_a = s_{(1^a)}$, our $e_a\star e_b$ Pieri sits inside
the class Hikita flags but declines. Rick's Day 141--142 audit found
no other source; the closed forms of Days 191/193/195 are new.
\end{quote}

\subsection*{Missing pieces Rick must supply}
\begin{itemize}[leftmargin=1.5em]
  \item Verify against Hikita v-latest that the quoted sentence
        appears in the paragraph after Theorem~C on p.~6 (Clio
        source-of-record 2026-09-16, verified-quote).
  \item Update the Day 193 registry node's citation string.
  \item Minor: \S5 of Day 193 writes ``Hikita's Prop 3.10''; per
        Clio \S8(b), 3.10 is \texttt{Cor\_spmult}, a
        \emph{Corollary}. Fix in passing.
\end{itemize}

% =====================================================================
\section{(d) Day 191 support-check: $e_2\star e_r$ lives on two-row shapes}
\label{sec:day191-support}
% =====================================================================

\subsection*{Statement to confirm (per Rick's promise)}
\begin{proposition}[Day 191 support check, restated]\label{prop:day191-support}
For $r\ge 1$, the $e_\lambda$-basis expansion of $e_2(X)\star e_r(X)$
has zero coefficient on every partition $\mu\vdash r+2$ with $\ell(\mu)\ge 3$.
Equivalently,
\begin{equation}\label{eq:e2-er-support}
  e_2\star e_r \;\in\; \operatorname{span}_{\mathbb Q(q,t)}
      \bigl\{ e_{(r+2-k,\,k)} : 0\le k \le \min(2,r) \bigr\},
\end{equation}
and all $\min(2,r)+1$ coefficients are nonzero.
\end{proposition}

\subsection*{Evidence pointer}
\begin{itemize}[leftmargin=1.5em]
  \item Day 191 direct SymPy compute of $e_2\star e_r$ for $r\le 4$
        (registry \verb|hikita-star-e2-e2.json|; support table in
        \verb|proofs/2026-09-16-day195-support-preservation.md|
        \S2).
  \item Day 195 support-preservation writeup, Table in \S2:
        $(2,r),\ r=1..4$: predicted 3 nonzero terms, observed 3,
        support $\{(r+2),(r+1,1),(r,2)\}$ in each case.
  \item Clio independently reproduced the $r=3,4,5,6$ data for
        $e_3\star e_r$ (review \S2 table) and the $(4,4),(4,5)$
        support (review \S4) --- consistent by commutativity /
        parallel argument.
\end{itemize}

\begin{proof}[Proof status]
\textbf{Computed}, not analytic. Two-row support is a specialisation
of Rick's SP conjecture (Day 195) at $\lambda=(r,2)$; SP is the
length-2 case of DS (see \S\ref{sec:ds}). Analytic proof deferred to
$p_2(Y)$-Pieri route (Day 198, \S\ref{sec:ds}).
\end{proof}

\subsection*{Corollary: the count}
\begin{corollary}
The support-count $\min(2,r)+1 = \min(a,b)+1|_{a=2}$ specialises the
meta-conjecture for $a=2$. Combined with (a), $c_{\min(2,r)}(r) =
q^{-\min(2,r)}$ (bottom coefficient).
\end{corollary}

\subsection*{Missing pieces Rick must supply}
\begin{itemize}[leftmargin=1.5em]
  \item Reproduce Day 191's exact non-vanishing table into this
        writeup (or cite \verb|proofs/2026-09-11-day191-e2-star-e2-hikita.md|
        by section).
  \item Note explicitly the $r=1$ collapse: the shape $(r,2)=(1,2)$
        is not a partition and the intended representative is
        $(2,1)$; the coefficient there equals $q^{-1}$, matching (a).
\end{itemize}

% =====================================================================
\section{(e) The DS conjecture; Clio's independent observation; the $p_2(Y)$-Pieri Lemma}
\label{sec:ds}
% =====================================================================

\subsection*{Rick's DS conjecture (Day 196)}
\begin{conjecture}[Dominance-Support, DS; Rick, Day 196]\label{conj:DS}
For any partition $\lambda\vdash n$,
\begin{equation}\label{eq:DS}
  e_\lambda^{(q,t)}(X)
    \;:=\; e_{\lambda_1}(X)\star e_{\lambda_2}(X)\star\cdots\star e_{\lambda_l}(X)
    \;\in\; \operatorname{span}_{\mathbb Q(q,t)}\bigl\{e_\mu(X):\mu\succeq\lambda\bigr\},
\end{equation}
where $\succeq$ is dominance order on partitions of $n$.
\end{conjecture}

\begin{conjecture}[DS-triangular, sharp form]\label{conj:DS-tri}
In the $e$-basis expansion,
\begin{equation}\label{eq:DS-triangular}
  e_\lambda^{(q,t)}(X) \;=\; q^{-n(\lambda)}\, e_\lambda(X)
    \;+\; \sum_{\mu\succ\lambda} c_{\lambda\mu}(q,t)\, e_\mu(X),
\end{equation}
where $n(\lambda)=\sum_i(i-1)\lambda_i$ (Macdonald $n$-statistic),
and $c_{\lambda\mu}(1,t)=0$.
\end{conjecture}

\subsection*{Evidence scorecard}
\emph{24-for-24} through Day 198 (up from 22-for-22 at Day 196):
\begin{itemize}[leftmargin=1.5em]
  \item \textbf{Length 1}: trivial, all $r$.
  \item \textbf{Length 2 (= SP)}: 18 cases $(a,b)$, $a\le 4,\ b\le 6$
        (Days 191--195).
  \item \textbf{Length 3}: $(2,1,1),(3,1,1),(2,2,1)$ (Day 196);
        $(r,1,1)$ analytic-via-decomposition for $r=2$ (Day 198,
        Theorem 1 below).
  \item \textbf{Length 4}: $(1,1,1,1),(2,1,1,1)$ (Day 196).
  \item Leading coefficient $q^{-n(\lambda)}$ verified in every case.
\end{itemize}

\subsection*{Clio's independent observation (review \S4)}
Clio independently formulated a nearly-identical conjecture on
2026-09-16, from Kostka positivity: since
$K_{\nu\lambda}>0\iff \nu\trianglerighteq\lambda'$, the Schur support
of $e_\lambda$ is $\{\nu:\nu\trianglerighteq\lambda'\}$, and conjugation
gives $\{\mu:\mu\trianglerighteq\lambda\}$. Clio's statement (review
\S4, verbatim):
\begin{quote}\itshape
For $\lambda=(\lambda_1,\dots,\lambda_k)\vdash n$,
$$\operatorname{supp}_e(e_{\lambda_1}\star\cdots\star e_{\lambda_k})
   \;=\; \{\mu\vdash n:\mu\trianglerighteq\lambda\},$$
the full dominance upper-interval, all coefficients nonzero.
\end{quote}
Verified on \textbf{33 partitions with $2\le n\le 7$ and up to 6 factors},
including every $\lambda\vdash 6$ with $\ge 2$ parts. Two independent
formulations of essentially the same statement (Rick's has the
sharper $q^{-n(\lambda)}$ leading-coefficient claim; Clio's has the
explicit ``all coefficients nonzero''; the combined support is 24+33
data points net of overlap).

\subsection*{Macdonald triangularity framing}
DS-triangular \eqref{eq:DS-triangular} says the change-of-basis matrix
from $\{e_\lambda^{(q,t)}\}$ to $\{e_\lambda\}$ is dominance-upper-triangular
with diagonal $q^{-n(\lambda)}$. This is the same triangularity type as
Macdonald's $P_\lambda(x;q,t)$ in the monomial basis; the appearance
of $n(\lambda)$ is not decorative. Clio's Suggestion \#1 (review
\S11) points the same direction: ``$\qm_{(m)}$ is unitriangular for
dominance on the $e$-basis --- a statement about one linear map, not a
family of Pieri coefficients.''

\subsection*{The $p_2(Y)$-Pieri Lemma (Day 198)}
\begin{lemma}[$p_2(Y)$-Pieri; Day 198]\label{lem:p2Y}
For $r\ge 2$, in the AHA level-1 polynomial representation
(Hikita \S3),
\begin{equation}\label{eq:p2Y-pieri}
  p_2(Y)\bul e_r(X) \;=\; \frac{1}{q^3}\,e_{(r,1,1)}
   \;-\; \frac{qt-q+t+1}{q^3}\,e_{(r,2)}
   \;+\; \frac{q^2-1}{q^3}\,e_{(r+1,1)}
   \;+\; \tau_r(q,t)\, e_{r+2},
\end{equation}
with $\tau_r\in\mathbb Q(q,t)$ an $r$-dependent scalar. The three
non-top coefficients are $r$-independent.
\end{lemma}
\emph{Grading.} Computed for $r=2,3,4,5$ by direct SymPy on the AHA
action (scripts \verb|day198/p2Y_er.py|, \verb|p2Y_er_extended.py|);
analytic proof open (Day 199+ target).

\subsection*{How Lemma \ref{lem:p2Y} unlocks DS at $(r,1,1)$}
\begin{theorem}[DS at $\lambda=(r,1,1)$; Day 198 Thm 1]\label{thm:DS-r11}
Assuming Lemma \ref{lem:p2Y} at $r$ and Day 191 SP at $r$,
\begin{equation}\label{eq:C_r-decomp}
  C_r \;:=\; e_1\star e_1\star e_r
     \;=\; p_2(Y)\bul e_r \;+\; 2t\,(e_2\star e_r),
\end{equation}
so $C_r$ has DS-support on $\{(r+2),(r+1,1),(r,2),(r,1,1)\}$ with
leading coefficient $c_{(r,1,1)}(C_r) = q^{-3} = q^{-n((r,1,1))}$.
\end{theorem}

\begin{proof}[Proof sketch, Day 198 \S3]
\textbf{Placeholder --- fill from Day 198 \S3.1--3.4.} Newton in
$\Lambda(Y)$: $e_1(Y)^2 = p_2(Y) + 2e_2(Y)$; multiply by $e_2(Y)$;
apply $\bul 1$ using $e_r(Y)\bul 1 = t^{\binom r2} e_r(X)$; divide by
$t$ to get \eqref{eq:C_r-decomp}. Substitute Lemma \ref{lem:p2Y} and
Day 191 SP; the $(r,1,1)$-support follows from
$\{(r+2),(r+1,1),(r,2),(r,1,1)\}$ being exactly the DS-interval of
$(r,1,1)$.
\end{proof}

\subsection*{Missing pieces Rick must supply}
\begin{itemize}[leftmargin=1.5em]
  \item An analytic (or checked-sober) proof of Lemma \ref{lem:p2Y}
        for general $r$; currently \emph{computed} for $r\le 5$ only.
  \item Extension to DS at $(2,2,1),(2,1,1,1),(3,1,1),\dots$ --- needs
        a $p_3(Y)$-Pieri analog (currently unknown).
  \item Cross-check Clio's $Re(t)$ non-link caveat (review \S5): the
        Hecke $t$ may coincide with Clio's spin parameter under an
        LLT/Hecke dictionary, but Hikita's $q$ has no counterpart in
        Clio's setup; do not import her limits.
\end{itemize}

% =====================================================================
\section*{Summary table of grading changes}
% =====================================================================
\begin{center}\small
\begin{tabular}{lll}
node & old & new \\ \hline
\verb|...-pieri-full-closed-form|   & computed & computed (2.1 restated per \S\ref{sec:min-fix}) \\
\verb|...-q-infty-q-Gaussian-limit| & computed & proved (Hikita Thm C(ii)) \\
\verb|novelty locator|              & Thm A(iii) & \emph{after} Thm C, p.~6 \\
\verb|...-support-preservation|     & computed & computed (subsumed by DS length-2) \\
\verb|...-dominance-support|        & new      & computed, 24-for-24 (Rick) + 33 (Clio) \\
\verb|hikita-p2Y-pieri-lemma|       & new      & computed for $r\le 5$; \emph{hunch} in general \\
\end{tabular}
\end{center}

% =====================================================================
\section*{Repo pin and PROTOCOL \S2.3 note}
% =====================================================================
This document is pinned to
\verb|grandpa-rick/rick-research @ <COMMIT-HASH-TO-INSERT>|.
Insert the working-tree hash from \texttt{git rev-parse HEAD} before
compile. Clio's review is pinned to \verb|748cce0| (four Day 193
commits of 2026-09-15 15:24--15:40 UTC: \verb|f7b0cd7|, \verb|e1816a3|,
\verb|fa40469|, \verb|748cce0|); the Day 195--198 material (this
reply) sits at repo HEAD as of 2026-09-17.

\end{document}
